% ============================================================
%  APÉNDICES
% ============================================================

\chapter{Manual de instalación y despliegue}
\label{ap:instalacion}

\section{Requisitos previos}
\label{ap:requisitos}

\begin{itemize}
  \item Docker Desktop~4.x o superior.
  \item Node.js~22~LTS y npm~10 (solo para desarrollo \textit{frontend} sin Docker).
  \item Python~3.12 (solo para desarrollo \textit{backend} sin Docker).
  \item Git.
\end{itemize}

\section{Despliegue con Docker Compose}
\label{ap:docker}

\begin{enumerate}
  \item Clonar los repositorios:
\begin{verbatim}
git clone <url-uniback>  Back/uniback
git clone <url-frontend> Front/gui-components
\end{verbatim}
  \item Copiar el fichero de variables de entorno y ajustar los valores:
\begin{verbatim}
cp .env.example .env
\end{verbatim}
  \item Levantar todos los servicios:
\begin{verbatim}
docker compose up --build
\end{verbatim}
  \item Abrir el navegador en \texttt{http://localhost:4200}.
\end{enumerate}

\section{Desarrollo local sin Docker}
\label{ap:dev-local}

\subsection{Backend}

\begin{verbatim}
cd Back/uniback
python -m venv .venv && source .venv/bin/activate
pip install -e ".[dev]"
uvicorn uniback.main:app --reload --port 8000
\end{verbatim}

\subsection{Frontend}

\begin{verbatim}
cd Front/gui-components
npm install
npx ng build ngt-gui          # compilar la librería
npx ng serve ngt-gui-app      # arrancar la aplicación
\end{verbatim}

% -------------------------------------------------------

\chapter{Referencia de extensiones x-* del esquema}
\label{ap:xextensions}

La tabla~\ref{tab:xextensions} recoge todas las extensiones \texttt{x-*} definidas por el
\textit{framework} y reconocidas por \textit{ngt-gui}.

\begin{longtable}{lp{4cm}p{6cm}}
  \caption{Extensiones x-* del esquema.} \label{tab:xextensions} \\
  \toprule
  \textbf{Extensión} & \textbf{Tipo} & \textbf{Descripción} \\
  \midrule
  \endfirsthead
  \multicolumn{3}{c}{\textit{(continuación)}} \\
  \toprule
  \textbf{Extensión} & \textbf{Tipo} & \textbf{Descripción} \\
  \midrule
  \endhead
  \bottomrule
  \endfoot
  \texttt{x-label} & \texttt{string} & Etiqueta legible para el campo en la \gls{ui}. \\
  \texttt{x-widget} & \texttt{string} & Tipo de widget Formly: \texttt{input},
    \texttt{select}, \texttt{textarea}, \texttt{date}, \texttt{checkbox}, \ldots \\
  \texttt{x-required} & \texttt{boolean} & El campo es obligatorio. \\
  \texttt{x-readonly} & \texttt{boolean} & El campo no es editable por el usuario. \\
  \texttt{x-order} & \texttt{integer} & Posición del campo en el formulario/tabla. \\
  \texttt{x-foreign-key} & \texttt{string} & Nombre de la entidad referenciada. \\
  \texttt{x-options-endpoint} & \texttt{string} & URL del endpoint para cargar opciones
    del selector. \\
  \texttt{x-validators} & \texttt{object} & Validadores adicionales: \texttt{maxLength},
    \texttt{minLength}, \texttt{pattern}, \texttt{min}, \texttt{max}. \\
  \texttt{x-hidden} & \texttt{boolean} & El campo no se muestra en la \gls{ui}. \\
  \texttt{x-default} & \texttt{any} & Valor por defecto para formularios de creación. \\
\end{longtable}

% -------------------------------------------------------

\chapter{Estructura del payload de la API REST}
\label{ap:api}

Todos los endpoints de \textit{uniback} envuelven sus respuestas en el objeto
\texttt{ResponseEnvelope}:

\begin{lstlisting}[language=json, caption={Estructura de ResponseEnvelope.}]
{
  "content": <datos | null>,
  "count":   <número total para listas | null>,
  "issues":  [{ "level": "warning|error", "message": "..." }]
}
\end{lstlisting}

Las peticiones de lista aceptan los parámetros de consulta:

\begin{table}[H]
  \centering
  \caption{Parámetros de consulta comunes a los endpoints de lista.}
  \label{tab:params-lista}
  \begin{tabular}{llp{7cm}}
    \toprule
    \textbf{Parámetro} & \textbf{Tipo} & \textbf{Descripción} \\
    \midrule
    \texttt{page}      & integer & Número de página (base 1). \\
    \texttt{page\_size} & integer & Registros por página (máx. 500). \\
    \texttt{search}    & string  & Búsqueda de texto completo. \\
    \texttt{order}     & string  & Campo y dirección (\texttt{name:asc}). \\
    \texttt{filter}    & JSON    & Filtros avanzados por campo. \\
    \bottomrule
  \end{tabular}
\end{table}
